
\section{Results and Evaluation}

We test our method on 13 indoor and outdoor scenes from three different standard datasets.
We evaluate on all scenes of Mip-NeRF 360 \cite{mipnerf360},
two scenes from Tanks \& Temples \cite{tnt},
as well as two scenes from Deep Blending \cite{db}.
Fig.~\ref{fig:default-eval} shows that our method faithfully reconstructs the input images.
%\Fredo{Should we discuss how we normalize and keep comparison fair wrt number of primitives vs memory size?}
% \begin{itemize}
% \item Fig. \ref{fig:results} line showing our result and GT on 2 scenes
% \item Show that quality is good
% \item Give Speed, memory etc.
% \end{itemize}

\subsection{Evaluation}
\label{sec:eval}

We compare our method against 2DGS \cite{2dgs} as a baseline,
and two 2DGS texturing methods,
namely (unpublished) BBSplat~\cite{bbsplat} and GStex \cite{gstex},
for which the code was available at the time of submission.
\revision{rev9584 2866 comparison to textured-3dgs 3dgs-mcmc}
{
Directly comparing to 3DGS-based models is not straightforward and might lead to misleading conclusions,
since these models have different properties and consistently perform better for NVS but worse in geometry reconstruction,
compared to 2DGS-based solutions~\cite{2dgs}.
For the same reason,
and because the code was not available at the time of submission,
we do not compare against \cite{textured3dgs},
which is a 3DGS-based method with textured primitives.
This model uses a pretrained 3DGS-MCMC model as initialization,
similar to GStex,
and includes RGBA textures,
similar to BBSplat.
We thus expect it to suffer from the disadvantages of both these two models (see discussion of the second experiment).
}
We ran these methods on all scenes using our machines to ensure fair comparison.
For 2DGS,
the normal consistency and depth distortion regularization terms were disabled,
as they are designed to enhance the geometric reconstruction
but often result in lower novel view synthesis (NVS) quality.

\begin{figure*}[!h]
\centering
\setlength{\tabcolsep}{1.5pt}

    \newcommand{\resultfigwidth}{.24\textwidth}

\begin{tabular}{cc|ccc}
    & Ground Truth & BBSplat & GStex & Ours \\
    \raisebox{0.4\height}{\rotatebox{90}{\small Mip-Nerf360}} &
    \spyimage{red}{\resultfigwidth}{images/table2/m360_gt.jpg}{1.4, 2.5}{1.1, 0.7} &
    \spyimage{red}{\resultfigwidth}{images/table2/m360_bbsplat_same_params_.jpg}{1.4, 2.5}{1.1, 0.7} &
    \spyimage{red}{\resultfigwidth}{images/table2/m360_gstex_output_same_params.jpg}{1.4, 2.5}{1.1, 0.7} &
    \spyimage{red}{\resultfigwidth}{images/table2/m360_texture-proj_final_new_downscale_threshold_0.02.jpg}{1.4, 2.5}{1.1, 0.7} \\
    
    \raisebox{0.3\height}{\rotatebox{90}{\small Deep Blending}} &
    \spyimage{red}{\resultfigwidth}{images/table2/db_gt.jpg}{2.5, 0.3}{4, 0.7} &
    \spyimage{red}{\resultfigwidth}{images/table2/db_bbsplat_same_params_.jpg}{2.5, 0.3}{4, 0.7} &
    \spyimage{red}{\resultfigwidth}{images/table2/db_gstex_output_same_params.jpg}{2.5, 0.3}{4, 0.7} &
    \spyimage{red}{\resultfigwidth}{images/table2/db_texture-proj_final_new_downscale_threshold_0.02.jpg}{2.5, 0.3}{4, 0.7} \\

    \raisebox{0.15\height}{\rotatebox{90}{\small Tanks\&Temples}} &
    \spyimage{red}{\resultfigwidth}{images/table2/tnt_gt.jpg}{2.1, 0.1}{1.1, 0.7} &
    \spyimage{red}{\resultfigwidth}{images/table2/tnt_bbsplat_same_params_.jpg}{2.1, 0.1}{1.1, 0.7} &
    \spyimage{red}{\resultfigwidth}{images/table2/tnt_gstex_output_same_params.jpg}{2.1, 0.1}{1.1, 0.7} &
    \spyimage{red}{\resultfigwidth}{images/table2/tnt_texture-proj_final_new_downscale_threshold_0.02.jpg}{2.1, 0.1}{1.1, 0.7} \\
\end{tabular}

\caption{
\label{fig:same-param-eval}
Adjusting BBSplat and GSTex to use the same budget as our method, results in significant visual degradation.}
\end{figure*}

\textbf{Qualitative Evaluation.}
%\Fredo{Sentence hard to parse. Simplify. Check I didn't break it. Is the end of the sentence useful?} 
In Figure~\ref{fig:default-eval},
we show visual results for the different models. We show
test views (i.e., not used for training) from one scene of each dataset.
Our method succeeds in reconstructing high-frequency details with high fidelity,
even with fewer parameters.
Similarly, Figure~\ref{fig:same-param-eval} displays visual results from the second experiment,
which highlight that other texturing methods struggle when constrained to the same parameter budget as ours.
GSTex uses a point cloud that is not trained with textured primitives in mind
and a static heuristic distribution of texels, which causes it to
not reconstruct well large parts of the scene.
BBSplat succeeds in geometrically reconstructing the scene, but exhibits texture stretching, which appears as blurriness.
This is due to the choice of mapping between intersection points and UV coordinates. 
%\Fredo{Hard to parse, long winded, rephrase}

%\subsection{Evaluation}
\textbf{Quantitative Evaluation.}
For our quantitative comparisons,
we present results for three error metrics, SSIM, PSNR, and LPIPS \cite{lpips},
that are typically used to evaluate NVS.
Additionally,
we report the number of primitives and
texels used,
which provides a precise measure of the resources required for each method
and demonstrates whether and how much textures contribute to the scene reconstruction.
\revision{rev9584 scope of the paper, primitive vs parameters}{
Finally,
since the models' capacity is directly tied to the number of parameters used,
we include the relevant column so that we can make fair comparisons and draw meaningful conclusions.
}

\revision{rev4160}{
The formula for the number of parameters differs slightly among the models.
}
For 2DGS, GSTex each primitive has 58 parameters (3 for position, 2 for scale, 4 for rotation, 1 for opacity, 48 for color),
while for BBSplat this number is 57,
as it lacks opacity.
In our method,
in addition to the parameters of 2DGS,
we include one additional per primitive for the texel size,
pushing the number to 59 parameters.
Regarding the parameters coming from texels,
GSTex and our method have 3 per texel,
while BBSplat has one additional,
corresponding to the alpha channel.


\begin{table*}[!ht]
\footnotesize
\setlength\tabcolsep{2pt}
\centering
\begin{tabular}{@{}>{\raggedright\arraybackslash}p{1.3cm}|>{\centering\arraybackslash}p{0.8cm}>{\centering\arraybackslash}p{0.8cm}>{\centering\arraybackslash}p{0.8cm}>{\raggedleft\arraybackslash}p{0.7cm}>{\centering\arraybackslash}p{0.9cm}>{\centering\arraybackslash}p{0.6cm}|>{\centering\arraybackslash}p{0.8cm}>{\centering\arraybackslash}p{0.8cm}>{\centering\arraybackslash}p{0.8cm}>{\raggedleft\arraybackslash}p{0.7cm}>{\centering\arraybackslash}p{0.9cm}>{\centering\arraybackslash}p{0.6cm}|>{\centering\arraybackslash}p{0.8cm}>{\centering\arraybackslash}p{0.8cm}>{\centering\arraybackslash}p{0.8cm}>{\raggedleft\arraybackslash}p{0.7cm}>{\centering\arraybackslash}p{0.9cm}>{\centering\arraybackslash}p{0.6cm}}
 & \multicolumn{6}{c|}{DeepBlending} & \multicolumn{6}{c|}{Mip-Nerf-360} & \multicolumn{6}{c}{Tanks\&Temples} \\
 & SSIM$\uparrow$ & PSNR$\uparrow$ & LPIPS$\downarrow$ & Texels & Params & FPS & SSIM$\uparrow$ & PSNR$\uparrow$ & LPIPS$\downarrow$ & Texels & Params & FPS & SSIM$\uparrow$ & PSNR$\uparrow$ & LPIPS$\downarrow$ & Texels & Params & FPS \\
\midrule
Points40k & 0.890 & 28.78 & 0.346 & 13.2M & 41.9M & 121 & 0.758 & 25.80 & 0.304 & 42.2M & 128.9M & 97 & 0.809 & 22.74 & 0.246 & 15.4M & 48.4M & 162 \\
Points80k & {\cellcolor{tabthird}} 0.898 & {\cellcolor{tabthird}} 29.56 & {\cellcolor{tabthird}} 0.328 & 15.5M & 51.2M & 88 & {\cellcolor{tabthird}} 0.777 & {\cellcolor{tabthird}} 26.39 & {\cellcolor{tabthird}} 0.285 & 43.9M & 136.5M & 79 & {\cellcolor{tabthird}} 0.823 & {\cellcolor{tabthird}} 23.25 & {\cellcolor{tabthird}} 0.235 & 12.3M & 41.7M & 108 \\
Points120k & {\cellcolor{tabsecond}} 0.903 & {\cellcolor{tabsecond}} 29.88 & {\cellcolor{tabsecond}} 0.315 & 17.9M & 60.7M & 75 & {\cellcolor{tabsecond}} 0.785 & {\cellcolor{tabsecond}} 26.68 & {\cellcolor{tabsecond}} 0.275 & 45.0M & 142.0M & 68 & {\cellcolor{tabsecond}} 0.830 & {\cellcolor{tabsecond}} 23.34 & {\cellcolor{tabsecond}} 0.232 & 10.9M & 39.6M & 96 \\
Points160k & {\cellcolor{tabfirst}} 0.905 & {\cellcolor{tabfirst}} 29.98 & {\cellcolor{tabfirst}} 0.308 & 19.9M & 69.1M & 56 & {\cellcolor{tabfirst}} 0.791 & {\cellcolor{tabfirst}} 26.88 & {\cellcolor{tabfirst}} 0.269 & 45.7M & 146.6M & 46 & {\cellcolor{tabfirst}} 0.833 & {\cellcolor{tabfirst}} 23.44 & {\cellcolor{tabfirst}} 0.227 & 10.5M & 40.8M & 88 \\
\end{tabular}

\caption{\label{tab:point-budget}
We run our technique with a limited primitive budget.
An increased primitive count increases quality,
without proportionally increasing the number of texels and parameters.
}
\end{table*}



In Table~\ref{tab:default-eval},
we compare methods trained with default settings.
In most cases,
our approach achieves competitive or superior quality across all three metrics,
while requiring a lower number of trainable parameters. %\TODO{Except for MipNerf -- say something}
Focusing on the composition of these parameters,
our converged models have significantly fewer,
but more expressive primitives than the baseline as illustrated in Fig.~\ref{fig:teaser} (left).
This highlights the ability of our model to account for the difference in complexity between geometry and appearance.
In contrast, GSTex uses a pretrained 2DGS point cloud and has a fixed budget of texels, distributed at initialization time; this results in limited usage of texture. 
This is in part because it starts with a converged set of primitives from 2DGS that is not well adapted to a solution with texture.
On the contrary, we optimize both primitives and textures from scratch.

Moreover,
while our primitive count is close to BBSplat's,
the total number of parameters used is significantly lower than theirs.
This can be attributed to our content-aware texel size determination,
which dynamically allocates model capacity to regions according to their needs,
as opposed to the fixed texel to primitive ratio that BBSplat imposes.
%\Fredo{Check that vs. which in English. }

To emphasize the importance of our content-aware texturing approach,
we next compare the other texturing methods by fixing both the primitive and texel budgets to ours.
In GSTex this is feasible,
because it allows the distribution of a given texel budget over a pretrained 2DGS point cloud.
We first trained 2DGS with the specified primitive budget and then gave the texel budget as input to GSTex.
However,
in BBSplat, controlling both the primitive and texel budgets is not possible,
as the method imposes a fixed ratio between the two,
determined by the texture resolution (16x16).
Therefore,
we calculated the number of primitives that would result in the same number of trainable parameters.
Table~\ref{tab:same-param-eval},
reports these results under these fixed parameters.
Note that BBSplat uses a skybox to model background and distant objects in outdoor scenes,
which was not taken into consideration when calculating the number of primitives.
This accounts for the slight discrepancy in the number of parameters.
Both the other methods observe a noticeable to significant drop in performance.
This is possibly due to the fact that model capacity is not distributed efficiently both across primitives and across parts of the scene.
In this setup, our method performs consistently better on average than previous solutions.
%\Fredo{Conclusion of the experiment?}

As a third experiment,
we run our method with a low, fixed primitive budget,
by skipping the splitting procedure whenever the model exceeds it.
The adaptive texel size strategies were left unchanged.
Table~\ref{tab:point-budget} shows the results.
As expected,
the performance gets better with an increasing primitive count,
as the use of primitives with a fixed gaussian falloff prevents faithfully reconstructing sharp edges with sparser point clouds.
However,
we note that the number of texels does not grow proportionally to the number of primitives
or even diminishes, in the case of Tanks~\&~Temples.
This is a direct result of our choice of texture coordinate mapping function and the adaptive texel size determination strategy.
Our separation of geometric and appearance parameters allows our models to automatically achieve a balance between the two that is appropriate for the scene.
\TODO{add BBSplat to show the proportional rise of parameters. Add the matplot graph}

% \begin{itemize}
% \item Compare to (3DGS), 2DGS* (no regul), GSTex, BBSplat, Ours 
% \item Methodology: left out, parameters for others etc.
% \end{itemize}


% \begin{itemize}
% \item On average same or better quality numbers and much lower number of parameters
% \item M360: low number of points, lower quality; discuss this
% \item However, BBsplat in this scene...
% \end{itemize}



% \begin{itemize}
% \item We are amazing with same memory
% \item explain how we did this for BBsplat, GSTex
% \end{itemize}
